\section{Discussion}
\label{sec:discussion}

\subsection{Why the local first-order conditions are insufficient}

The accepted uniform-price vector is not generated by an algebraic mistake in the smooth first-order conditions. It is the correct intersection of those conditions on the lower switching branch. The problem is global. Demand contains a flat no-switch region. For the smaller firm, moving from the smooth poaching optimum to the right-hand plateau boundary can sharply increase price without sacrificing its inherited customers. The profitability of that finite deviation is exactly what \(x_H\) measures.

This distinction matters methodologically. In piecewise price games, verifying first-order and second-order conditions on one regime does not establish Nash equilibrium unless all feasible regime changes and kinks are checked. Here, the omitted comparison is economically meaningful: it is the choice between continuing to poach the dominant firm's customers and ceasing poaching while harvesting one's own inherited base.

\subsection{Why the welfare correction is not simply the equilibrium correction}

Two logically separate issues arise downstream. First, the uniform-price vector is an equilibrium only on the corrected domain. Second, even as a counterfactual price profile, its allocation changes at \(x_0=x_u\). Extending a one-way-switch consumer-surplus integral below that cutoff reverses the order of an integration interval and no longer represents realized consumer choices.

Consequently, one cannot repair the accepted analysis by adding the condition \(x_0\ge x_H\) alone. Consumer surplus must also be integrated over the actual clipped allocation. This is why the weak consumer-surplus reversal disappears even at the profile level, while the welfare ranking favoring uniform pricing survives on the equilibrium overlap.

\subsection{Economic implications inside the model}

The corrected overlap between weak HBP and a pure uniform-price equilibrium is small: normalized switching cost must satisfy \(s<s_c\approx0.1334\). In that region, HBP raises consumer surplus but lowers the smaller firm's profit and lowers total welfare relative to uniform pricing. The welfare effect reflects real switching and transport costs, not price transfers.

Under strong inherited dominance, the source consumer-surplus comparison is more sensitive to price-domain conventions because below-cost poaching prices can support the same zero-sales allocation. Social welfare is more robust: the entire zero-sales price family induces the same allocation, so the strong welfare ranking is selection-invariant.

\subsection{Limitations}

The analysis does not solve the mixed-strategy uniform-price game below \(x_H\). It therefore does not compare equilibrium consumer surplus or welfare in that region. Nor does it claim robustness to alternative spatial demand, outside options, heterogeneous switching costs, or imperfect recognition. Those are distinct models rather than omitted cases of the theorem.

A final source limitation is textual rather than mathematical. Equation-level statements are verified against a post-referee accepted manuscript. Until the exact typeset Version of Record is compared directly, this paper does not claim that every equation number or algebraic display in the final publisher layout is identical to that accepted version.
